%%% SECTION 9: A CANCER PATIENT'S JOURNEY THROUGH PHYSICAL AI %%%

A Cancer Patient's Journey, Physical AI \cite{patient-journey-paper} provides
definitive evidence that Physical AI can manage a complete oncology trial
patient journey autonomously across ten stages, with documented safety
outcomes, regulatory compliance, and cost savings. This single-patient
simulation documents a Stage IIIB non-small cell lung cancer (NSCLC) patient
over 1,120 days, demonstrating end-to-end Physical AI capability from
prescreening through closeout. An example is illustrated in \autoref{Motion}. Together with the 24-hour multi-patient
simulation \cite{site-docs}, this constitutes pivotal evidence that
state-of-the-art AI has the ability to understand, work with, and scale
Physical AI applications beyond human and prior technology capabilities.

\subsection{Ten-Stage Pipeline Architecture}
\label{subsec:journey-pipeline}

The patient journey is organized into ten sequential stages, each documented
with regulatory mapping, clinical activities, Physical AI system involvement,
and outcomes. The stages span the complete trial lifecycle, demonstrating that
Physical AI is not limited to isolated robotic procedures but can orchestrate
the entire patient experience.

\begin{table}[H]
\centering
\caption{Ten-Stage Physical AI Patient Journey Pipeline}
\label{tab:journey-stages}
\small
\begin{tabularx}{\textwidth}{c l X}
\toprule
\textbf{Stage} & \textbf{Name} & \textbf{Physical AI Activities} \\
\midrule
1 & Prescreening & AI-assisted eligibility assessment, record review \\
2 & Screening & Physical AI diagnostic imaging, biomarker analysis \\
3 & Informed Consent & AI-enhanced consent with robot demonstrations \\
4 & Enrollment & Automated registration, baseline data collection \\
5 & Treatment Planning & Digital twin modeling, treatment optimization \\
6 & Surgical Intervention & Robot-assisted tumor resection \\
7 & Post-Surgical Care & Robotic monitoring, rehabilitation support \\
8 & Follow-up Treatment & AI-guided therapy, companion robot support \\
9 & Long-term Monitoring & Automated surveillance, imaging follow-up \\
10 & Closeout & Automated reporting, regulatory documentation \\
\bottomrule
\end{tabularx}
\end{table}

The ten-stage pipeline demonstrates end-to-end capability that differentiates
Physical AI oncology trials from prior robotic surgery studies focusing only on
surgical outcomes. Each stage is documented with three-perspective ASCII
diagrams showing the regulatory, clinical, and patient views simultaneously.

\subsection{Pre-Trial Stages: Prescreening Through Enrollment}
\label{subsec:journey-pretrial}

Stages 1 through 4 cover the pre-trial process from initial patient
identification through formal enrollment. Stage 1 (Prescreening) uses AI
algorithms to review patient medical records against trial eligibility
criteria, identifying candidates with Stage IIIB NSCLC who meet inclusion
criteria and have no disqualifying conditions. Physical AI systems assist by
automating record retrieval, cross-referencing imaging databases, and flagging
potential candidates for physician review.

Stage 2 (Screening) deploys Physical AI imaging systems for diagnostic
confirmation, including robotic CT-guided biopsy for tissue sampling and
AI-assisted pathology analysis for tumor characterization. Stage 3 (Informed
Consent) implements the Physical AI-specific consent requirements from the
adapted 21 CFR Part 50 \cite{cfr50-adapt}, including robot demonstrations,
AI transparency explanations, and documentation of patient rights under
AB 2847 from the site documentation package \cite{site-docs}.

Stage 4 (Enrollment) uses automated systems for patient registration, baseline
data collection, and treatment arm assignment. The regulatory framework mapping
shows how each pre-trial stage complies with the adapted standards: the
Adaption: ICH Harmonised Guideline \cite{ich-adapt} governs investigator
conduct at each stage, the adapted 21 CFR Part 50 governs informed consent and
human subjects protection, and the adapted 21 CFR Part 312 \cite{cfr312-adapt}
governs IND compliance and Phase 0 validation completion.

\subsection{Treatment Stages: Surgical and Therapeutic Interventions}
\label{subsec:journey-treatment}

Stages 5 through 8 cover the active treatment period. Stage 5 (Treatment
Planning) uses digital twin technology to create a patient-specific
computational model of the tumor and surrounding anatomy, enabling simulation
of multiple treatment approaches before selecting the optimal strategy. AI
algorithms analyze the digital twin against historical outcomes data to predict
treatment response.

Stage 6 (Surgical Intervention) documents the robot-assisted tumor resection
performed by a surgical robot system. The procedure is monitored through the
Physical AI safety matrix requirements, with continuous force, position, and
vital sign monitoring. Surgical outcomes include procedure duration, blood loss,
margin clearance, and immediate post-operative status. The adverse event
profile documents any robot-related events during the procedure.

Stage 7 (Post-Surgical Care) deploys rehabilitation exoskeletons for mobility
support, companion robots for emotional support, and monitoring cobots for
continuous vital sign surveillance. Stage 8 (Follow-up Treatment) covers
adjuvant therapy with AI-guided dosing optimization and Physical AI-assisted
treatment delivery.

USL scoring \cite{usl-standard} is applied during treatment stages to verify
that each robot involved in patient care meets minimum readiness thresholds.
Surgical robots must achieve Advanced band (USL 7.0 or higher) for Dimension D
before performing procedures, while support robots (cobots, companions) may
operate at Intermediate band (USL 5.0 or higher).

\subsection{Post-Treatment Stages: Closeout and Analysis}
\label{subsec:journey-closeout}

Stages 9 and 10 cover the post-treatment period through trial closeout. Stage 9
(Long-term Monitoring) uses automated Physical AI surveillance systems for
scheduled imaging, blood work monitoring, and symptom assessment over the
1,120-day follow-up period. AI algorithms analyze trends in monitoring data to
detect potential recurrence or late treatment effects. Stage 10 (Closeout)
automates the generation of regulatory documentation including case report
forms, safety narratives, and protocol compliance reports, ensuring that all
Physical AI-generated data is properly archived per the retention requirements
established in the adapted standards.

\subsection{Regulatory Compliance Coverage}
\label{subsec:journey-compliance}

The patient journey demonstrates compliance with all three adapted regulatory
standards at every stage. The regulatory framework mapping documents which
specific provisions of each adapted standard apply at each stage, creating a
traceability matrix from clinical activities to regulatory requirements.

\begin{table}[H]
\centering
\caption{Regulatory Standard Coverage Across Journey Stages}
\label{tab:journey-compliance}
\small
\begin{tabularx}{\textwidth}{c X X X}
\toprule
\textbf{Stage} & \textbf{ICH E6(R3)} & \textbf{21 CFR Part 50} & \textbf{21 CFR Part 312} \\
\midrule
1-2 & Investigator qualifications & Subject identification & IND eligibility \\
3-4 & Informed consent procedures & Consent elements & Enrollment procedures \\
5-6 & Quality management & Safety matrix & Phase reporting \\
7-8 & Safety reporting & Ongoing consent & Adverse event reporting \\
9-10 & Record keeping & Data protection & Annual/final reports \\
\bottomrule
\end{tabularx}
\end{table}

\subsection{Physical AI Ecosystem and Data Flow}
\label{subsec:journey-ecosystem}

Data flows between Physical AI systems, clinical staff, regulatory systems, and
patients throughout the ten-stage journey. The Physical AI ecosystem comprises
multiple robot types operating in coordinated workflows, with data exchange
managed through MCP servers \cite{national-mcp-paper} using the Model Context
Protocol \cite{mcp-protocol}. Each robot generates data streams that are
captured, validated, and routed to appropriate destinations: clinical databases
for care coordination, regulatory systems for compliance reporting, and
research databases for outcome analysis.

The data flow architecture ensures that no single robot operates in isolation.
Surgical robots share procedure data with monitoring cobots, which in turn
update patient records accessed by companion robots for post-procedure support.
This integrated data flow supports the continuous monitoring requirements of
the adapted 21 CFR Part 50 \cite{cfr50-adapt} and the data governance
provisions of the adapted ICH E6(R3) \cite{ich-adapt}.

\subsection{FDA Cost-Savings Analysis}
\label{subsec:journey-cost}

The patient journey includes a detailed financial analysis documenting
per-patient cost reductions at each stage. Physical AI reduces costs through
automated screening (Stage 1-2), streamlined consent processes (Stage 3),
efficient enrollment (Stage 4), AI-optimized treatment planning (Stage 5),
precision surgical intervention reducing complications and hospital stay
(Stage 6), automated post-surgical monitoring reducing nursing workload
(Stage 7), AI-guided follow-up reducing unnecessary clinic visits (Stage 8-9),
and automated closeout documentation reducing administrative burden (Stage 10).

Tufts CSDD data \cite{tufts-delay} establishes that each day of delay in drug
development has significant economic consequences. The timeline compression
demonstrated in the patient journey, from faster screening and enrollment
through more efficient treatment delivery and automated documentation,
translates directly to accelerated drug development timelines. NCI data on
trial costs \cite{nci-trials-paying} provides context for understanding the
economic burden that Physical AI can reduce.

\subsection{Trial Timeline Comparison}
\label{subsec:journey-timeline}

The Physical AI trial timeline introduces the Phase 0 simulation validation
stage but achieves net time savings through compression at every subsequent
phase. Enrollment speed increases through automated screening and AI-assisted
eligibility verification. Treatment throughput increases through 24-hour
robotic operations and reduced per-procedure times. Data processing accelerates
through automated documentation and real-time regulatory reporting through MCP
servers. The net effect is a significantly shorter overall trial timeline
despite the addition of the Phase 0 requirement.

\subsection{Discussion of Patient Journey Evidence}
\label{subsec:journey-discussion}

The patient journey evidence is significant for several reasons. First, it
demonstrates autonomous workflow capability across all ten stages, establishing
that Physical AI is not limited to individual procedures but can manage
complete trial operations. Second, it validates the regulatory adaptations by
showing that a real patient scenario can be mapped to the adapted standards at
every stage. Third, it provides the financial data that underpins the economic
analysis in Section~\ref{sec:financial}. Fourth, it establishes safety
performance benchmarks against which future Physical AI trials can be measured.

Both the single-patient journey and the 24-hour multi-patient simulation are
pivotal completed examples demonstrating state-of-the-art AI ability to
understand, proficiently work with, and scale Physical AI applications. The
evidence is comprehensive: safety data, regulatory compliance, cost analysis,
and patient outcomes all support the transition to Physical AI oncology trials.

\subsection{Limitations and Future Directions}
\label{subsec:journey-limitations}

The single-patient simulation has acknowledged limitations: it documents a
simulated rather than real patient, covers only one cancer type (Stage IIIB
NSCLC), and represents outcomes from AI-generated scenarios rather than
clinical observations. These limitations are partially addressed by the
24-hour multi-patient simulation covering 15 cancer types and 168 patients.
Full validation requires real-world clinical trials, which this National
Platform is designed to enable through its comprehensive regulatory, standards,
infrastructure, and documentation framework.

\subsection{Detailed Stage Timeline and Physical AI Systems}
\label{subsec:journey-detailed-timeline}

\begin{longtable}{c p{2.25cm} p{6cm} p{6cm}}
\caption{Detailed Stage Timeline with Physical AI Systems Involved}
\label{tab:journey-timeline} \\
\toprule
\textbf{Stage} & \textbf{Duration} & \textbf{Physical AI Systems} & \textbf{Key Deliverables} \\
\midrule
\endfirsthead
\toprule
\textbf{Stage} & \textbf{Duration} & \textbf{Physical AI Systems} & \textbf{Key Deliverables} \\
\midrule
\endhead
\bottomrule
\endfoot
1 & Days 1-14 & AI screening algorithms, imaging robots & Eligibility determination, baseline imaging \\
2 & Days 15-28 & Needle-placement robot, cobots & Biopsy results, biomarker analysis, staging confirmation \\
3 & Days 29-35 & Companion robots, AI consent system & Signed Physical AI-specific consent, patient education \\
4 & Days 36-42 & Cobots, data management AI & Enrollment confirmed, baseline data collection complete \\
5 & Days 43-70 & Digital twin AI, treatment planning AI & Personalized treatment plan, simulation validation \\
6 & Days 71-85 & Surgical robot (da Vinci), monitoring cobots & Tumor resection, procedure documentation \\
7 & Days 86-120 & Rehab exoskeleton, companion robots, cobots & Post-surgical recovery milestones, wound healing \\
8 & Day 121-365 & Cobots (drug prep), imaging robots, AI monitoring & Adjuvant therapy delivery, response assessment \\
9 & 366-1050 & Imaging robots, AI surveillance, companion robots & Long-term monitoring, recurrence screening \\
10 & 1051-1120 & AI documentation, data management systems & Final analysis, regulatory closeout, data archive \\
\end{longtable}

The total journey spans 1,120 days (approximately 3 years), consistent with
typical oncology trial timelines for Stage IIIB NSCLC. Physical AI system
involvement is continuous from Day 1 through Day 1,120, with different robot
categories taking primary roles at different stages. The surgical robot is
active for a concentrated period (Days 71-85), while companion robots provide
support across nearly all stages. Imaging robots are involved in prescreening
(Stage 1), treatment monitoring (Stages 7-9), and surveillance (Stage 9),
demonstrating their utility across the complete trial lifecycle.

\subsection{Cost-Savings Detail by Stage}
\label{subsec:journey-cost-detail}

\begin{table}[H]
\centering
\caption{Physical AI Cost Impact by Journey Stage}
\label{tab:journey-cost-stage}
\small
\begin{tabularx}{\textwidth}{c l X}
\toprule
\textbf{Stage} & \textbf{Primary Savings Mechanism} & \textbf{Impact} \\
\midrule
1-2 & Automated screening & Reduced coordinator time per patient, faster eligibility determination \\
3 & AI-enhanced consent & Streamlined consent process with automated documentation \\
4 & Automated enrollment & Reduced data entry time, fewer transcription errors \\
5 & Digital twin planning & Reduced planning iterations, optimized treatment strategy first time \\
6 & Robotic precision & Shorter procedure duration, reduced complication rate, shorter hospital stay \\
7 & Automated monitoring & Reduced nursing surveillance hours, earlier complication detection \\
8 & AI-guided therapy & Optimized dosing reducing waste, fewer dose modifications needed \\
9 & Automated surveillance & Reduced follow-up visits through remote monitoring, faster imaging analysis \\
10 & Automated closeout & Reduced administrative burden, faster regulatory submission \\
\bottomrule
\end{tabularx}
\end{table}

The cumulative effect of cost savings across all ten stages compounds as the
number of patients increases. At the single-patient level, each stage
contributes incremental savings. At the 168-patient level demonstrated in the
24-hour simulation \cite{site-docs}, the per-stage savings multiply across the
patient cohort while fixed costs (infrastructure, robot procurement) remain
constant, creating the favorable scale economics analyzed in
Section~\ref{sec:financial}.

\subsection{Safety Performance Across Stages}
\label{subsec:journey-safety}

The patient journey documents safety performance at each stage using metrics
aligned with the adapted ICH E6(R3) \cite{ich-adapt} adverse event
classification. Stage 6 (Surgical Intervention) carries the highest inherent
risk due to direct Physical AI involvement in an invasive procedure, with the
surgical robot classified as Tier 3 under the adapted 21 CFR Part 50
\cite{cfr50-adapt} classification system. The safety matrix was completed
before surgery, runtime monitoring tracked force, position, and vital signs
continuously, and post-procedure documentation captured all surgical
parameters.

Stages 1-4 and 9-10 carry lower risk profiles (primarily Tier 1 and Tier 2
Physical AI interactions) but still require documentation per the adapted
standards. The companion robot interactions throughout the journey represent
Tier 1 activities with standard monitoring. The overall safety record across
all ten stages supports the conclusion that Physical AI can manage the
complete trial lifecycle without compromising patient safety.

\subsection{Comparison with Traditional Oncology Trial Timelines}
\label{subsec:journey-traditional-comparison}

Traditional oncology trials for Stage IIIB NSCLC typically require 4-6 months
for enrollment (compared to 42 days in the Physical AI journey), 6-12 months
for active treatment (compared to approximately 11 months), and 2-3 years for
follow-up (comparable to the Physical AI timeline). The most significant
timeline compression occurs in the pre-treatment stages (screening, consent,
enrollment, treatment planning) where Physical AI automation reduces weeks of
manual coordination to days of AI-assisted processing. Treatment stage
compression comes primarily from 24-hour robot availability rather than
shift-limited human scheduling, and from AI-optimized treatment planning
reducing the number of plan iterations.

The Phase 0 simulation validation required by the adapted 21 CFR Part 312
\cite{cfr312-adapt} adds time before patient contact. However, Phase 0
validation can proceed in parallel with traditional preclinical and regulatory
preparation, and the thorough pre-validation it provides may reduce the time
needed for early clinical phases by demonstrating safety upfront rather than
through gradual dose escalation.

\subsection{Expanded Cost-Savings Analysis by Stage}
\label{subsec:journey-cost-expanded}

The cost-savings mechanisms summarized in Table~\ref{tab:journey-cost-stage}
can be quantified more precisely by examining the specific labor, material, and
time savings at each stage. The following analysis draws on data from the
single-patient journey simulation \cite{patient-journey-paper}, the 24-hour
multi-patient simulation \cite{site-docs}, and contextual cost data from Tufts
CSDD \cite{tufts-delay} and NCI \cite{nci-trials-paying}.

\subsubsection{Pre-Trial Stage Savings (Stages 1-4)}

Stages 1 through 4 represent the patient acquisition pipeline where Physical
AI achieves the most dramatic time compression relative to traditional trials.
In traditional oncology trials, the pre-trial period is the most
labor-intensive per-patient phase, requiring clinical research coordinators to
manually review medical records, coordinate screening appointments, guide
patients through consent, and complete enrollment documentation.

Stage 1 (Prescreening) deploys AI algorithms to review patient records
against trial eligibility criteria. Traditional prescreening requires an
average of 4-6 coordinator hours per candidate to review medical history,
imaging records, pathology reports, and medication lists. AI-assisted
prescreening reduces this to approximately 30 minutes of coordinator time per
candidate for review and confirmation of AI-flagged results. For a trial
targeting 168 patients (the scale of the 24-hour simulation), this represents
a savings of approximately 588-924 coordinator hours during the prescreening
phase alone. At prevailing coordinator compensation rates, this translates to
significant per-patient cost reduction before any patient enters the screening
phase.

Stage 2 (Screening) uses Physical AI imaging systems for diagnostic
confirmation. Robotic CT-guided biopsy reduces procedure setup time through
automated patient positioning and equipment calibration. AI-assisted pathology
analysis accelerates tissue characterization from days to hours for
preliminary results. The combined effect reduces the screening period from a
traditional 3-4 weeks to approximately 2 weeks, with each day of acceleration
contributing to faster enrollment and reduced trial timeline.

Stage 3 (Informed Consent) benefits from AI-enhanced consent processes that
provide standardized, comprehensive information delivery while maintaining the
human interaction essential for valid consent. The Physical AI consent process
includes robot demonstrations that give patients direct experience with the
systems they will encounter, which research suggests improves comprehension and
reduces consent-related delays from patient uncertainty. Documentation
automation reduces the administrative burden of consent processing from
approximately 2-3 hours per patient to approximately 45 minutes.

Stage 4 (Enrollment) uses automated systems for registration and baseline data
collection that reduce transcription errors and data entry time. Traditional
enrollment requires manual entry of patient demographics, medical history,
baseline assessments, and treatment arm assignments into multiple systems.
Automated enrollment with AI-verified data entry reduces this to a
single-entry process with automated cross-system propagation, saving
approximately 1-2 hours of coordinator time per patient while reducing data
errors that can require costly corrections later in the trial.

\subsubsection{Treatment Stage Savings (Stages 5-8)}

Treatment stages generate cost savings through precision, efficiency, and
continuous operation that are not achievable with human-only workflows.

Stage 5 (Treatment Planning) uses digital twin technology to simulate multiple
treatment approaches before selecting the optimal strategy. Traditional
treatment planning requires multiple iterations of physician review,
dosimetrist calculation, physics quality assurance, and plan revision, a
process that typically spans 2-4 weeks. Digital twin-based planning with AI
optimization reduces this to approximately 1-2 weeks by generating and
evaluating candidate plans computationally before physician review, presenting
the clinical team with pre-optimized options rather than requiring iterative
manual refinement.

Stage 6 (Surgical Intervention) achieves cost savings through precision that
reduces complications and hospital stay. Published data on robotic surgical
outcomes in oncology demonstrate reduced blood loss, shorter hospital stays,
and faster return to normal activity compared to open surgery. In the
simulation, robotic surgical intervention achieved complete tumor resection
with clear margins in a single procedure, avoiding the need for revision
surgery that occurs in a percentage of traditional cases. Each avoided
complication represents substantial cost savings in additional treatment,
extended hospitalization, and delayed progression to follow-up therapy.

Stage 7 (Post-Surgical Care) deploys automated monitoring that reduces nursing
surveillance hours while improving detection of post-surgical complications.
Continuous robotic vital sign monitoring with AI-analyzed trend detection
provides earlier warning of complications compared to periodic manual checks.
Rehabilitation exoskeletons accelerate mobility recovery, potentially reducing
the number of physical therapy sessions required. Companion robots provide
continuous patient support that supplements but does not replace human nursing
care, enabling nursing staff to focus on higher-acuity tasks.

Stage 8 (Follow-up Treatment) benefits from AI-guided therapy optimization
that reduces drug waste and minimizes dose modifications. AI algorithms
analyzing real-time patient response data can identify the optimal dosing
strategy faster than traditional empirical adjustment, reducing the number of
dose modification cycles and associated monitoring visits. Automated drug
preparation by cobots reduces pharmacist preparation time and virtually
eliminates compounding errors.

\subsubsection{Post-Treatment Stage Savings (Stages 9-10)}

Stage 9 (Long-term Monitoring) generates the largest per-patient savings
through reduction of unnecessary follow-up visits. Traditional long-term
monitoring requires scheduled in-person visits at regular intervals regardless
of patient status. AI-powered surveillance can stratify patients by recurrence
risk and tailor monitoring intensity accordingly, reducing low-risk patient
visits while maintaining or increasing monitoring for high-risk patients. For
the Stage IIIB NSCLC patient documented in the journey, the monitoring period
spans approximately 685 days (Days 366-1050). Even a modest reduction in the
number of in-person visits during this period, replaced by remote AI-analyzed
monitoring, generates substantial cumulative savings across the patient cohort.

Stage 10 (Closeout) uses automated documentation to reduce the administrative
burden that often delays trial completion. Traditional closeout requires manual
compilation of case report forms, safety narratives, protocol compliance
reports, and data queries. Automated closeout with AI-generated documentation
reduces this from weeks of administrator time to days, with the AI system
compiling comprehensive regulatory documentation from the structured data
collected throughout the trial. Faster closeout directly contributes to faster
regulatory submission, which translates to faster drug approval timelines.

\subsubsection{Scale Economics}

The per-patient savings at each stage compound as the trial scales from single
patients to full enrollment. Fixed costs (Physical AI system procurement, site
infrastructure, staff training) are amortized across larger patient cohorts,
while per-patient variable costs decrease through automation and efficiency.
The 24-hour simulation treating 168 patients with 29 robots at 99.7 percent
uptime demonstrated that the Physical AI infrastructure can support high
throughput without proportional increases in variable costs. The marginal cost
of treating each additional patient through the Physical AI pipeline is lower
than the marginal cost in a traditional trial because robotic systems operate
continuously without shift-change delays, overtime costs, or fatigue-related
efficiency losses.

Tufts CSDD data \cite{tufts-delay} establishes that each day of delay in drug
development carries significant economic consequences for sponsors. The
timeline compression demonstrated across the ten stages, particularly in
pre-trial acceleration (Stages 1-4) and closeout automation (Stage 10),
translates directly to reduced time-to-market for investigational drugs. When
multiplied across the portfolio of oncology drugs in development, the
aggregate economic impact of Physical AI-enabled timeline compression is
substantial, providing a compelling economic case for Physical AI adoption
that complements the clinical quality and patient safety arguments.
